% Generated draft: requires mathematical and source review.
\documentclass[11pt]{article}
\usepackage[margin=25mm]{geometry}
\usepackage{fontspec}
\setmainfont{Latin Modern Roman}
\usepackage{amsmath,amssymb,mathtools,unicode-math}
\setmathfont{Latin Modern Math}
\usepackage{enumitem,needspace,xurl,hyperref}
\setlength{\parindent}{0pt}
\setlength{\parskip}{5pt}
\setlength{\emergencystretch}{3em}
\newcommand{\N}{\mathbb{N}}
\newcommand{\Z}{\mathbb{Z}}
\newcommand{\Q}{\mathbb{Q}}
\newcommand{\R}{\mathbb{R}}
\newcommand{\C}{\mathbb{C}}
\newcommand{\F}{\mathbb{F}}
\newcommand{\A}{\mathbb{A}}
\newcommand{\PP}{\mathbb{P}}
\newcommand{\E}{\mathbb{E}}
\newcommand{\eps}{\varepsilon}
\newcommand{\dd}{\,\mathrm{d}}
\begin{document}
\section*{UnsolvedMath ID: 11100016}
\textbf{Source problem code:} catalogue entry
\subsection{Definitions and mathematical statement}
Fix a prime \(p\), a height-\(n\) formal group \(\mathbb G\) over the intended perfect field \(k\) of characteristic \(p\), and let \(\mathbb S=\operatorname{Aut}_k(\mathbb G)\) be its profinite stabilizer group. In a usual Lubin--Tate coordinate model the universal deformation ring is
\[R\cong W(k)[[u_1,\ldots,u_{n-1}]],\]
and Morava \(E\)-theory is even periodic with \(\pi_0E\cong R\) and \(\pi_*E\cong R[u,u^{-1}]\), \(|u|=-2\). Exact field, coordinate and extended-stabilizer conventions are not specified in the item.

The source already reports a stabilizer action from the Hopkins--Miller theorem. Its main request is instead a construction of \(E\) making that action manifest from the deformation-moduli input: a functorial construction on the formal-group/deformation groupoid, whose automorphisms induce compatible self-equivalences of the resulting spectrum, with a continuous \(\mathbb S\)-action. The expected action on \(R\) is the continuous action on Lubin--Tate deformations; coherent spectral lifts and an appropriate profinite continuity model must be supplied. Continuity only after forgetting topology or only on an abstract automorphism group is not the requested explanation.

Hovey describes the historical proof as calculating an \(A_\infty\)-automorphism space of \(E\), finding it homotopy discrete with abstract \(\pi_0\) isomorphic to \(\mathbb S\), and asks for a more natural method. He suggests passing from Lubin--Tate moduli to spectra through some infinite-loop-space technology, explicitly as a presumed approach, not an exclusive required mechanism. His hope that the proof would be published is a dated source remark.

The research summary reports a modern solution via Goerss--Hopkins/Lurie and an \(E_\infty\)-refinement. A spectral moduli-sheaf interpretation would use sections of an appropriate sheaf of spectra over the Lubin--Tate deformation space \(\operatorname{Spf}R\). Ordinary sections of an ordinary structure sheaf are just a ring and do not themselves specify \(E\); the local deformation space is also not the entire global moduli stack of formal groups. The summary's conflated notation and truncated formulation require this distinction before treating them as a precise construction theorem. The source does not fix a particular continuous-action formalism, so none is silently chosen.
\subsection{Short English statement}
Construct Morava \(E\)-theory functorially from Lubin--Tate data so its stabilizer action, including profinite continuity, is built in. The record reports a modern solution, while its spectral-sheaf and bibliographic details need clarification.
\subsection{Sources}
\begin{itemize}
\item Local UnsolvedMath registry, permanent ID \(11100016\), exact source code \texttt{catalogue entry}.
\item Mark Hovey, \emph{Hovey's Algebraic Topology Problem List}, Morava \(K\)- and \(E\)-theory, item 6, surviving mirror: \url{https://www-users.cse.umn.edu/~tlawson/hovey/}. Registry access 2026-07-29; html-ordered-list extraction.
\item Hopkins--Miller are credited by the source for the known action. The summary names Goerss--Hopkins, \emph{Moduli spaces of commutative ring spectra}, a claimed 2014 publication context, Goerss's \emph{Hopf algebroids and the structure of MU-modules}, and Lurie, but its entries/publication description are mixed and its final text is truncated.
\end{itemize}
\end{document}
